% !TEX root = ../../../tree_notes.tex
% =====================================================================
%  Chapter 11 opener -- "Chapter at a Glance", three reading paths + the
%  ladder staircase, the two running examples (the word set
%  WORDS = {car, cares, care, cat, dog} and the XOR set A = [3,10,5,25,2,8]),
%  the "who owns what" comparison (trie vs BST vs hash map vs segment/
%  Fenwick-on-prefix), roadmap, recognition guide, constraint->tool table,
%  notation.
%  NO numbered boxes here. Exports fig:c08-wordstrie and fig:c08-triebasics.
% =====================================================================

\begin{chaptertheorybox}{Chapter at a Glance}
Every tree so far has compared \emph{values} stored at nodes: the BSTs of Chapters~5--6
(\chref{c05}, \chref{c06}) decide a left/right turn by $x<$ or $x>$ the node; the heaps of
Chapter~10 (\chref{c07}) decide only ``does a parent beat a child.'' This chapter inverts the idea.
A \term{trie} (from \emph{re}\textbf{trie}\textbf{val}) is a tree whose \textbf{EDGES carry the
information}: one character (or one bit) per level, so a key is read by \emph{walking}, symbol by
symbol, from the root. The payoff is a single reframed cost. Searching a BST for a length-$L$ string
costs $\bigO{L\log N}$ -- $\log N$ node visits, each doing up to $L$ characters of comparison. A trie
costs $\bigO{L}$ \emph{no matter how many keys $N$ you stored}: the number of levels is the key length,
not the item count (\thmref{c08-searchcost}).
\begin{itemize}[nosep]
  \item \textbf{One symbol per level.} An edge is labelled by a character (or a bit). Membership,
        prefix queries, and ``enumerate everything under this prefix'' all cost $\bigO{L}$, independent
        of $N$. \hfill \lv{2}--\lv{3}
  \item \textbf{The node is a small dispatch table.} A node holds the map (or array) ``symbol $\to$
        child pointer'' plus an \term{end-of-word} flag. Layout is a real design choice -- fixed
        \code{child[26]} vs a \code{map} vs a sorted edge list -- and it drives the memory budget.
        \hfill \lv{3}
  \item \textbf{Numbers are just short strings.} Insert a $\le\!10^9$ integer as its $30$ bits, MSB
        first, and the same machine answers \emph{maximum-XOR} questions greedily (\lemref{c08-greedybit}).
        \hfill \lv{4}
\end{itemize}

\textbf{You will learn to}
\begin{itemize}[nosep]
  \item build a string trie and write \code{insert} / \code{search} / \code{startsWith} cold; \hfill\lv{2}--\lv{3}
  \item turn a trie into a dictionary engine -- wildcard search, prefix counting, autocomplete top-$k$,
        and grid word search (LC~212); \hfill\lv{4}
  \item lay a set of integers out as a bitwise trie and read off the maximum XOR pair in $\bigO{L}$; \hfill\lv{4}
  \item compress single-child chains into a radix trie, and \emph{build a working persistent XOR-trie}
        that answers ``max XOR over a version or a value range''; \hfill\lv{5}
  \item recognise the trie signature in a problem statement and pick string-vs-binary, array-vs-map,
        plain-vs-persistent without hesitation. \hfill\lv{4}
\end{itemize}

\textbf{Prerequisites.} Trees, nodes, edges and DFS from Chapters~1--4; the BST as an ordered tree
(\chref{c05}) which we contrast against exactly once; a hash map's $\bigO{1}$-amortised lookup as the
rival structure. The size-$k$ \emph{heap} behind top-$k$ suggestions is \emph{cited} from
Chapter~10 (\chref{c07}), not re-derived. A little bit arithmetic ($x\oplus y$, shifting) is all 8.3
needs. Persistence mechanics in general live in Chapter~21 (\chref{c19}); this chapter builds one
concrete persistent trie and forwards the abstraction.
\end{chaptertheorybox}

% ---------------- Reading paths + ladder position ----------------
\begin{chaptertheorybox}{How to Read This Chapter (three paths)}
Level badges $\lv1$--$\lv6$: $\lv2$ is a bare insert/search fact, $\lv3$ is a solid LC-Medium trie
design you code straight from the definition, $\lv4$ is where you must \emph{combine} the trie with
another idea (a grid DFS, a bitwise greedy, a subtree count), and $\lv5$ is the Summit and the
persistent trie, where the difficulty is holding the right invariant.

\textbf{Fast path ($\lv3$, ``I just want a prefix tree'').} Read \secref{c08-triebasics} and the solved
Implement Trie (LC~208), then skim the applications you actually need (wildcard, autocomplete). Do the
practices LC~677, LC~648. You leave able to cold-solve ``store words, search a word, does a prefix
exist'' and any dictionary-shaped Medium.

\textbf{Structural path ($\lv3$--$\lv4$).} Add \secref{c08-applications} (wildcard DFS, prefix counting,
Word Search II) and \secref{c08-binary} (the bitwise trie and the max-XOR greedy). Order:
8.1 $\to$ 8.2 $\to$ 8.3. The bridge fact is that a bit is just a two-letter character -- once you see
that, the XOR greedy is the same walk you already know.

\textbf{Full path.} Everything in order, finishing at \secref{c08-variants} -- the radix/compressed trie,
the memory cost table, and a \emph{working} persistent XOR-trie (LC~1707) -- and the Summit
(CF~706D \emph{Vasiliy's Multiset}), the purest dynamic bitwise trie: add, erase, and query the maximum
XOR against the current set, all $\bigO{30}$, with deletes handled by a per-node count.

The forward rope is Chapter~21 (\chref{c19}) for general persistence and Chapters~13--14
(\chref{c13}, \chref{c14}) for value-domain Fenwick/segment k-th-with-deletes; the backward rope is
Chapter~10 (\chref{c07}) for the heap inside top-$k$, and Chapters~5--6 (\chref{c05}, \chref{c06}) for
the ordered tree the trie deliberately is not.
\end{chaptertheorybox}

\begin{figure}[H]
\centering
\begin{tikzpicture}[diagram, x=9mm, y=7mm]
  \foreach \st/\lab in {1/lv1,2/lv2,3/lv3,4/lv4,5/lv5,6/lv6} {
    \fill[\lab!22, draw=\lab, line width=0.8pt] (\st-1,0) rectangle (\st,\st*0.52);
    \node[font=\scriptsize\bfseries, text=\lab] at (\st-0.5,\st*0.52-0.42) {L\st};}
  \node[annot, anchor=south west, font=\scriptsize]
    at (0.05,0.62) {a tree that reads a key symbol by symbol};
  \node[annot, anchor=south west, font=\scriptsize]
    at (1.05,1.14) {8.1 layout; insert/search/startsWith};
  \node[annot, anchor=south west, font=\scriptsize]
    at (2.05,1.66) {8.2 wildcard DFS; prefix count};
  \node[annot, anchor=south west, font=\scriptsize]
    at (3.05,2.18) {8.2 Word Search II; 8.3 max-XOR greedy};
  \node[annot, anchor=south west, font=\scriptsize, text=lv5]
    at (4.05,2.70) {8.4 radix; persistent XOR-trie};
  \node[annot, anchor=south west, font=\scriptsize, text=lv5]
    at (5.05,3.22) {summit (CF 706D)};
\end{tikzpicture}
\caption*{Chapter~11 on the six-step ladder: the layout and the three string operations are
$\lv2$--$\lv3$; wildcard DFS, prefix counting and the max-XOR greedy sit at $\lv3$--$\lv4$; Word
Search II and the bitwise design climb to $\lv4$; the compressed and persistent tries reach $\lv5$,
where the Summit (CF~706D \emph{Vasiliy's Multiset}) turns the whole chapter into one add / erase /
query walk on a bitwise trie. Fast path $=$ steps 2--4; full path runs to 6.}
\end{figure}

% ---------------- The running examples ----------------
\subsection*{The Running Examples: One Word Set, One Number Set}

Two small inputs, reused in every section so the ideas always land on something concrete.

\textbf{Word set $\mathcal{W}$} $= \{\text{car},\ \text{cares},\ \text{care},\ \text{cat},\ \text{dog}\}$
-- five words, but note the shared prefixes: \emph{car}, \emph{care}, \emph{cares} all run down one
spine and fork only at their last letters, while \emph{cat} shares \emph{ca} then leaves, and \emph{dog}
is a disjoint branch. $\mathcal{W}$ is the running example for \code{insert}/\code{search}/
\code{startsWith} (\secref{c08-triebasics}), wildcard and prefix-counting walks
(\secref{c08-applications}), and autocomplete. Storing these five words costs \emph{nodes}, not words:
the trie has $13$ non-root nodes no matter that ``care'' is inside three of them.

\textbf{Number set $A$} $= [3, 10, 5, 25, 2, 8]$ -- the official LeetCode~421 input, whose answer is
$\boxed{28}$ (from $5\oplus 25$). Read each value as a $5$-bit string and the same symbol-by-symbol
walk, now choosing the \emph{opposite} bit whenever it exists, finds the maximum-XOR pair
(\secref{c08-binary}). $A$ is reused for the bitwise-trie greedy, and its ``prefer the opposite bit''
rule is exactly what the Summit (CF~706D) runs dynamically.

\begin{figure}[H]
\centering
\begin{tikzpicture}[diagram, scale=0.9, transform shape,
  level 1/.style={sibling distance=42mm},
  level 2/.style={sibling distance=26mm},
  level 3/.style={sibling distance=17mm},
  level 4/.style={sibling distance=13mm},
  level 5/.style={sibling distance=11mm},
  tend/.style={tnode, fill=successgreen!20, draw=successgreen, thick}]
  \node[troot] {$\cdot$}
    child { node[tnode] {c}
      child { node[tnode] {a}
        child { node[tnode] {r}
          child { node[tend] {e}
            child { node[tend] {s} } } }
        child { node[tend] {t} } } }
    child { node[tnode] {d}
      child { node[tnode] {o}
        child { node[tend] {g} } } };
  \node[annot, align=left, font=\scriptsize] at (-4.0,2.1)
    {\textbf{green $=$ end-of-word}\\ \emph{car} at $r$,\\ \emph{care} at $e$,\\ \emph{cares} at $s$};
  \node[annot, align=left, font=\scriptsize] at (4.2,0.4)
    {\emph{cat} leaves at $t$;\\ \emph{dog} is a separate\\ branch sharing no prefix};
\end{tikzpicture}
\caption{The trie of $\mathcal{W}=\{\text{car},\text{cares},\text{care},\text{cat},\text{dog}\}$. Each
edge is one character; the root is blank ($\cdot$). Only four nodes are green end-of-word markers -- the
word \emph{car} ends at the node $r$ that also continues into \emph{care}/\emph{cares}, which is exactly
why a ``search'' (must end on a green node) differs from a ``startsWith'' (may stop anywhere): checking
\texttt{care} walks $c$-$a$-$r$-$e$ and sees green, checking \texttt{car} walks $c$-$a$-$r$ and sees green
too, but \texttt{careful} would walk $c$-$a$-$r$-$e$ and then find no $f$ edge.}\label{fig:c08-wordstrie}
\end{figure}

\begin{figure}[H]
\centering
\begin{tikzpicture}[diagram,
  nd/.style={boxnode, fill=primaryblue!10, draw=primaryblue, minimum size=7mm, font=\small},
  edg/.style={-{Stealth[length=2mm]}, deepblue, thick}]
  \node[nd] (p) at (0,0) {$v$};
  \node[boxnode, fill=white, draw=darkgray, font=\scriptsize, text width=22mm, right=13mm of p] (child)
    {child[``a''] $\to w$\\ child[``c''] $\to$ null\\ \ldots\\ \textcolor{successgreen}{\texttt{end}=true}};
  \draw[edg] (p) -- node[annot, font=\scriptsize, above, align=center, yshift=1mm]
    {edge label $=$\\ one symbol} (child);
  \node[annot, align=left, font=\scriptsize] at (0,-1.55)
    {a node $=$ dispatch\\ table + end flag};
\end{tikzpicture}
\caption{Anatomy of one trie node: a small table mapping each outgoing symbol to a child pointer, plus a
single \term{end-of-word} boolean saying ``a stored key terminates exactly here.'' Nothing about the key
is stored \emph{in} the node -- the key is the \emph{path} from the root. Swap the table for a fixed
\code{child[26]} array, a \code{map}, or a \code{vector<pair<char,node>>} and everything else in this
chapter is unchanged (\defref{c08-nodelayout}).}\label{fig:c08-triebasics}
\end{figure}

\medskip
\noindent Four structures answer superficially similar ``do I have this key?'' questions; knowing which
owns what stops you reaching for the wrong engine:
\begin{center}
\renewcommand{\arraystretch}{1.25}
\setlength{\tabcolsep}{5pt}
\begin{tabular}{@{}p{0.20\linewidth} p{0.41\linewidth} p{0.30\linewidth}@{}}
\toprule
\textbf{Structure} & \textbf{What it is for} & \textbf{Where it lives} \\
\midrule
Trie / prefix tree & keys read \emph{symbol by symbol}: $\bigO{L}$ membership \emph{and} every
  prefix/range-of-keys query, independent of $N$ & this chapter (\secref{c08-triebasics}) \\
Hash map / set & $\bigO{1}$-amortised \emph{exact} membership -- but \emph{no} prefix queries, no order,
  no ``list keys starting with \texttt{ca}'' & not a tree; the foil in \secref{c08-triebasics} \\
BST / balanced tree & keys \emph{ordered}: predecessor, successor, any rank -- but $\bigO{L\log N}$ per
  string compare & Chapters~5--6 (\chref{c05}, \chref{c06}) \\
Segment / Fenwick on a value domain & range aggregates and $k$-th-with-counts over an \emph{array of
  positions}, not over a key's spelling & Chapters~13--14 (\chref{c13}, \chref{c14}) \\
\bottomrule
\end{tabular}
\end{center}

% ---------------- Roadmap diagram ----------------
\begin{figure}[H]
\centering
\begin{tikzpicture}[diagram, node distance=3.5mm,
    top/.style={boxnode, fill=primaryblue!10, draw=primaryblue, text width=22mm,
                minimum height=11mm, font=\small},
    res/.style={annot, text width=22mm, anchor=north, font=\scriptsize}]
  \node[top] (t1) {\textbf{8.1} String Trie};
  \node[top, right=of t1] (t2) {\textbf{8.2} Applications};
  \node[top, right=of t2] (t3) {\textbf{8.3} Binary Trie};
  \node[top, right=of t3] (t4) {\textbf{8.4} Variants};
  \foreach \a/\b in {t1/t2, t2/t3, t3/t4} \draw[flow] (\a) -- (\b);
  \node[res] at ([yshift=-2mm]t1.south) {layout; insert, search, startsWith};
  \node[res] at ([yshift=-2mm]t2.south) {wildcard; prefix count; Word Search II};
  \node[res] at ([yshift=-2mm]t3.south) {bits as edges; max-XOR greedy};
  \node[res] at ([yshift=-2mm]t4.south) {radix compress; persistent trie};
\end{tikzpicture}
\caption*{The four topics are one idea -- a key is its path from the root -- seen on four alphabets.
8.1 \emph{lays} the alphabet as letters and writes the three operations; 8.2 spends the walk on
dictionary questions (a wildcard is ``try every edge'', a prefix count is ``count this subtree'');
8.3 changes the alphabet to $\{0,1\}$ and the goal to a greedy (take the opposite bit); 8.4 engineers
the structure -- collapse chains (radix), keep past versions (persistent) -- which is what the Summit
and LC~1707 demand.}
\end{figure}

% ---------------- Topic table ----------------
\begin{center}
\small
\renewcommand{\arraystretch}{1.25}
\setlength{\tabcolsep}{4pt}
\begin{tabular}{@{}p{0.235\linewidth} p{0.365\linewidth} c p{0.23\linewidth}@{}}
\toprule
\textbf{Topic} & \textbf{Key results} & \textbf{Levels} & \textbf{Solved problem} \\
\midrule
\secref{c08-triebasics}\ \ The Trie of Strings
  & edge-per-symbol node, $\bigO{L}$ search, node layouts, insert/search/startsWith & $\lv2$--\lv{3}
  & \probref{c08-implement} \\
\addlinespace
\secref{c08-applications}\ \ Dictionary Applications
  & wildcard DFS, prefix/subtree count, autocomplete top-$k$, Word Search II & $\lv3$--\lv{4}
  & \probref{c08-wordsearch} \\
\addlinespace
\secref{c08-binary}\ \ The Binary (XOR) Trie
  & bits MSB-first, prefer-opposite-bit greedy, max-XOR pair & $\lv4$
  & \probref{c08-maxxor} \\
\addlinespace
\secref{c08-variants}\ \ Compressed \& Persistent
  & radix/single-child collapse, path-copying persistent XOR-trie & $\lv4$--\lv{5}
  & \probref{c08-persistxor} \\
\bottomrule
\end{tabular}
\end{center}

% ---------------- Problem recognition guide ----------------
\subsection*{Problem Recognition Guide --- When You See \emph{this}, Think \emph{that}}

Trie problems announce themselves with a \emph{key that has an internal structure you read a piece at a
time}. If the phrasing mentions ``words / dictionary / prefix / autocomplete'', or ``maximum XOR'', or
``many strings sharing beginnings'', that is the tell; every row points at the section that builds the
tool.

\begin{center}
\renewcommand{\arraystretch}{1.3}
\setlength{\tabcolsep}{4pt}
\begin{tabular}{@{}p{0.44\linewidth} p{0.34\linewidth} l@{}}
\toprule
\textbf{When the problem says\ldots} & \textbf{\textcolor{deepblue}{\textbf{Think\ldots}}} & \textbf{Section} \\
\midrule
"store a dictionary; does this word / this prefix exist?"
  & a string trie; walk letters, end-flag distinguishes search vs startsWith & \secref{c08-triebasics} \\
\addlinespace
"replace each word by its shortest prefix" / "count by prefix"
  & trie walk; the first end-node is the answer, or subtree counts & \secref{c08-applications} \\
\addlinespace
"'.' matches any letter" / "search a word with wildcards"
  & DFS branching over all child edges at a '.' & \secref{c08-applications} \\
\addlinespace
"find many words in a letter grid at once"
  & build a trie of the words, DFS the grid matching it & \secref{c08-applications} \\
\addlinespace
"maximum XOR of two numbers / with an element"
  & binary trie, prefer the opposite bit at the top & \secref{c08-binary} \\
\addlinespace
"max XOR where the value is $\le m$, or within a range/version, add/erase live"
  & persistent or offline binary trie with per-node counts & \secref{c08-variants} \\
\bottomrule
\end{tabular}
\end{center}

\begin{center}
\renewcommand{\arraystretch}{1.25}
\setlength{\tabcolsep}{4pt}
\begin{tabular}{@{}p{0.29\linewidth} p{0.42\linewidth} p{0.20\linewidth}@{}}
\toprule
\textbf{Constraint you see} & \textbf{What it forces} & \textbf{Chapter-11 tool} \\
\midrule
total inserted letters $\le 10^{5}$--$10^{6}$, many prefix queries & $\bigO{L}$ per query, one node per letter & string trie \\
small alphabet (a--z, or bits $\{0,1\}$) & fixed \code{child[26]} / \code{child[2]} arrays beat maps in speed & \secref{c08-triebasics} \\
huge alphabet or very sparse branching & \code{map} / \code{vector} edges, or a radix trie, to save memory & \secref{c08-variants} \\
value $\le 10^{9}$, "maximize XOR" & a $30$/$31$-level binary trie, greedy on bits & \secref{c08-binary} \\
also bound the chosen value by $m$ / a range / a version & persistent or offline trie with counts & \secref{c08-variants} \\
\bottomrule
\end{tabular}
\end{center}

\begin{notebox}[Not this chapter?]
If you only ever ask ``is this \emph{exact} key present?'' and never touch prefixes or order, a hash map
is simpler and faster -- a trie earns its keep precisely on \emph{prefix} and \emph{shared-beginning}
work. If you need a \emph{total order} on keys (predecessor, successor, arbitrary rank), that is the
ordered trees of Chapters~5--6 (\chref{c05}, \chref{c06}). If you need multi-pattern \emph{substring}
matching (``does any dictionary word appear anywhere in this text?''), that is Aho--Corasick -- a trie
plus failure links -- which is previewed in \secref{c08-variants} but \emph{not} built here. A trie owns
exactly one thing and owns it perfectly: \emph{work that reads a key symbol by symbol and never depends on
how many keys are stored}.
\end{notebox}

% ---------------- Notation ----------------
\subsection*{Notation Used in This Chapter}

\begin{center}
\renewcommand{\arraystretch}{1.2}
\setlength{\tabcolsep}{4pt}
\begin{tabular}{@{}p{0.20\linewidth} p{0.27\linewidth}@{\quad}p{0.17\linewidth} p{0.27\linewidth}@{}}
\toprule
\textbf{Symbol} & \textbf{Meaning} & \textbf{Symbol} & \textbf{Meaning} \\
\midrule
$L$                 & key length in symbols (chars, or bits) -- the trie's search cost & $N$ & number of stored keys (which the cost ignores) \\
$\mathcal{W}$        & the running word set $\{\text{car},\text{cares},\text{care},\text{cat},\text{dog}\}$ & $A$ & the number set $[3,10,5,25,2,8]$ (LC 421) \\
$\mathit{child}[c]$  & the edge from a node labelled by symbol $c$ & $\mathit{end}$ & end-of-word flag at a node \\
$x \oplus y$         & bitwise XOR; a $1$ bit where $x,y$ differ & $b_k\!\ldots b_0$ & bits of a value, read MSB-first \\
$C$                  & value bound (e.g.\ $2^{30}$ for $\le\!10^9$); trie depth $\lceil\log_2 C\rceil$ & $\bigstar$ & marks the Summit \\
root & the blank node all paths start from; holds no symbol & & \\
\bottomrule
\end{tabular}
\end{center}
